\documentclass[10pt,a4paper]{article}
\usepackage[margin=17mm]{geometry}
\usepackage[T1]{fontenc}
\usepackage[utf8]{inputenc}
\usepackage{lmodern,amsmath,amssymb,booktabs}
\setlength{\parindent}{0pt}
\setlength{\parskip}{3.5pt}
\pagestyle{empty}
\newcommand{\ext}{\mathrm{ext}}
\begin{document}

{\large\bfseries Extremal mass at fixed angular momentum in EGB, $D=5$}\\[1pt]
{\small Status as of 9 October 2026. Asymptotically flat black holes, two equal spins $J_1=J_2=J$, $G_5=1$, action $R+\alpha\mathcal L_{\rm GB}$.}

\medskip
We want the smallest mass $M_\ext(J,\alpha)$ that still has a horizon when
$J$ is held fixed and the Gauss--Bonnet coupling $\alpha$ is varied. By
dimensional analysis it is a function of one variable:
\[
M_\ext=\pi^{1/3}J^{2/3}\,m(s),\qquad s=\frac{\pi^{2/3}\alpha}{J^{2/3}} .
\]
No closed form for $m(s)$ is known. This is what we have.

\textbf{1. Weak coupling: the series through $s^4$.}
\[
m(s)=\frac32+s-\frac{236}{105}s^2+d_3\,s^3+d_4\,s^4+O(s^5),
\]
\[
d_3=\frac{60851752}{33075}-1536\,\zeta(3)=-6.548307\ldots,\qquad
d_4=-\frac{592541765354296}{1489863375}+\frac{227328}{5}\zeta(3)+\frac{88064}{25}\pi^4=66.100522\ldots
\]
The linear term is known [1]. The coefficients $d_2$ and $d_3$ were first
obtained from the Euclidean action. We have now rederived them in a
different way: we solve for the exterior metric of the extremal black hole
order by order in $\alpha$, exactly, and read the mass directly off the
fall-off of $g_{tt}$ at infinity. At each order the metric functions are
combinations of multiple polylogarithms in $1/r^2$ with rational
coefficients, of weight at most the order. The integration constants are
fixed by asymptotic flatness and by requiring the horizon to be a double
zero of $g^{rr}$ (extremality). This is how $d_4$ was obtained; before, it
was known only numerically. An independent numerical calculation (spectral
boundary-value problem) gives $66.10052$, which agrees in all seven
available digits.

Checks: at each order the solution satisfies the reduced field equations
exactly, and the horizon angular velocity, computed separately from the
mass, satisfies the extremal first law $dM_\ext=2\Omega_H\,dJ$ (at fixed
$\alpha$) with zero residual at orders 2, 3 and 4.

Caveat: $d_4$ is exact \emph{provided} certain constants that appear when the
polylogarithms are evaluated at the horizon equal the zeta values we
identified numerically (100 digits, error $<10^{-80}$). This is not proved.
The series is also formal: we do not know whether it converges. The
coefficients ($1,\,-2.2,\,-6.5,\,66$) do not yet show a clear growth pattern.

\textbf{2. Strong coupling ($\alpha\to\infty$, $J$ fixed): a candidate.}
A matched analysis of three regions (horizon throat, core, intermediate zone)
gives
\[
M_\ext=\frac{3\pi\alpha}{4}+\frac{4J}{3\sqrt\alpha}+o\!\left(J/\sqrt\alpha\right),\qquad \Omega_H\sqrt\alpha\to\tfrac23 ,
\]
consistent with the first law. It relies on the regions matching
uniformly, which is not proved. The full numerical solution reaches only
$\alpha\simeq3.12$ (with $J/\alpha^{3/2}\simeq0.26$), too far from the limit
to confirm it.

\textbf{3. The point $J=0$.} The term $3\pi\alpha/4$ is the mass of the
static solution with zero horizon radius, which is singular. Kleihaus, Kunz
and Radu conjecture that the rotating extremal branch ends there [2]. It is
not a regular extremal black hole: static black holes with a horizon always
have positive temperature.

\textbf{4. What is missing.} There is no closed form for $m(s)$ at finite
coupling, and no proof that none exists. The negative results we have apply
to specific families: Hamilton--Jacobi superpotentials with finitely many
powers of $r$, Liouvillian solutions of one perturbation equation (its
differential Galois group is $SL(2,\mathbb C)$), and a candidate
differential equation for $m(s)$ that an independent numerical point rules
out. The second-order exterior already contains $\mathrm{Li}_2(1/r^2)$, and
$\zeta(3)$ and $\pi^4$ appear at higher orders, which makes a simple
algebraic formula unlikely.

\textbf{Next steps.} The $s^5$ term with the same program (the constants it
needs are already identified); with five or six coefficients the radius of
convergence can be estimated. On the numerical side, extending the branch
beyond $\alpha\simeq3$ to test the strong-coupling limit.

\vfill
{\footnotesize [1] Ma, Li, L\"u, arXiv:2009.00015. \quad
[2] Kleihaus, Kunz, Radu, arXiv:2303.12471. \quad
Code and outputs: \texttt{extremal-mass/code/}; see \texttt{extremal-mass/README.md}.}
\end{document}
